\documentclass[11pt]{amsart}
\usepackage[margin=1.1in]{geometry}
\usepackage{amsmath,amssymb,amsthm,mathtools}
\usepackage{enumitem}
\usepackage{booktabs}
\usepackage{xcolor}
\usepackage[colorlinks=true,linkcolor=blue!50!black,citecolor=green!40!black,urlcolor=blue!60!black]{hyperref}

\newtheorem{theorem}{Theorem}[section]
\newtheorem{proposition}[theorem]{Proposition}
\newtheorem{lemma}[theorem]{Lemma}
\newtheorem{conjecture}[theorem]{Conjecture}
\theoremstyle{definition}
\newtheorem{definition}[theorem]{Definition}
\newtheorem{problem}[theorem]{Problem}
\newtheorem{heuristic}[theorem]{Heuristic}
\theoremstyle{remark}
\newtheorem{remark}[theorem]{Remark}
\newtheorem{example}[theorem]{Example}

\newcommand{\F}{\mathbb{F}}
\newcommand{\Z}{\mathbb{Z}}
\newcommand{\R}{\mathbb{R}}
\newcommand{\C}{\mathbb{C}}
\newcommand{\CFK}{\mathit{CFK}}
\newcommand{\HFK}{\mathit{HFK}}
\newcommand{\CF}{\mathit{CF}}
\newcommand{\HF}{\mathit{HF}}
\newcommand{\HM}{\mathit{HM}}
\newcommand{\Spinc}{\mathrm{Spin}^c}
\newcommand{\CSD}{\mathcal{L}}
\newcommand{\Tdual}{T^{\vee}}
\newcommand{\cB}{\mathcal{B}}
\newcommand{\cC}{\mathcal{C}}
\newcommand{\cG}{\mathcal{G}}
\newcommand{\cM}{\mathcal{M}}
\newcommand{\gr}{\mathrm{gr}}
\DeclareMathOperator{\Cone}{Cone}

\title[Towards a gauge-theoretic $\CFK^\infty$]{Towards a gauge-theoretic $\CFK^{\infty}$:\\
anatomy of the knot Floer complex and an Atiyah--Floer proposal}
\date{October 2026}

\begin{document}
\begin{abstract}
No instanton or monopole theory currently produces a complex over $\F[U,V]$, or an analogue of $\CFK^\infty$.
We propose one, built on Atiyah--Floer principles.

We first take the full knot Floer complex $\CFK^\infty$ apart into the mechanisms that produce it.  This yields a
specification: $\CFK_{\F[U,V]}$ is the Floer complex of the closed manifold $Y$, with $U$ given, plus an integer Alexander
function and a one-sided positivity of a $z$-charge.

On the gauge side, the base is the monopole complex (or the $SO(3)$-equivariant instanton complex) of $Y$, for a metric
stretched along a doubly pointed Heegaard surface, and trajectories are weighted by $V^{n_z}$.
\begin{itemize}
\item For monopoles, $n_z$ counts zeros of the vortex component of the spinor along the strip swept by the $z$-arc of $K$.
That component is the $E_-(\dot K)$-eigenline of Clifford multiplication by the tangent of $K$.
\item For instantons, $n_z$ is a relative index of the Dirac operator coupled to the instanton and to a unit $U(1)$ flux
string along the $z$-arc.  In the Atiyah--Floer limit it becomes intersection with a theta-type divisor in the moduli
space of the Heegaard surface, and the knot enters through how its determinant line is trivialised over the handlebody
Lagrangians.
\end{itemize}
Positivity is inherited from holomorphicity in the adiabatic limit, through an integrality bootstrap.

We describe a neck-free local model at critical Aharonov--Bohm holonomy $\beta=\tfrac12$.  There the two
$r^{-1/2}$ edge modes are the components seen by $w$ and $z$, the two arcs of $K$ meet at ``creases'', and the creases
become Zemke's dividing sets on cobordisms.  A separate singular variant is also treated.  Its two variables are the
nef ends of the segment of Kähler classes of parabolic moduli spaces (with $t=4\alpha$), and it is a Khovanov-side
theory, not $\CFK$.

We isolate the analytic problems: the adiabatic limit with a transverse knot, admissibility, edge calculus, a
failing $C^0$ bound and crease bubbles, and compatible divisors.  Genus-one splittings (pillowcase with two marked
points; $\mathrm{Sym}^1T^2$) are proposed as the first rigorous test.  Statements labelled Theorem, Proposition or Lemma
are proved here or standard; everything else is labelled Conjecture, Heuristic or Problem.
\end{abstract}
\maketitle
\setcounter{tocdepth}{1}
\tableofcontents

%======================================================================
\section{The object: $\CFK^\infty$ and $\CFK_{\F[U,V]}$}\label{sec:cfk}
%======================================================================

Throughout, $Y$ is a rational homology sphere, $K\subset Y$ is an oriented null-homologous knot,
$\F=\Z/2$, and $\mathfrak s\in\Spinc(Y)$ is suppressed when it plays no role.

\subsection{Definition}
A \emph{doubly pointed Heegaard diagram} $\mathcal H=(\Sigma,\boldsymbol\alpha,\boldsymbol\beta,w,z)$ for
$(Y,K)$ is a Heegaard diagram for $Y$ with two basepoints in $\Sigma\setminus(\boldsymbol\alpha\cup\boldsymbol\beta)$
such that $K$ is the union of an arc from $w$ to $z$ in $\Sigma\setminus\boldsymbol\alpha$ pushed into
$U_\alpha$ and an arc from $z$ to $w$ in $\Sigma\setminus\boldsymbol\beta$ pushed into $U_\beta$.
Equivalently (and this is the description that matters below): if $f$ is a self-indexing Morse
function with one minimum $p_0$ and one maximum $p_3$, then $K=\gamma_w\cup\overline{\gamma_z}$
where $\gamma_w,\gamma_z$ are the two \emph{upward} gradient flow lines from $p_0$ to $p_3$ through
$w$ and $z$.  The orientation of $K$ agrees with the flow on $\gamma_w$ and disagrees with it on $\gamma_z$.

\begin{definition}[\cite{OS-knots,Rasmussen}]
$\CFK_{\F[U,V]}(\mathcal H)$ is the free $\F[U,V]$-module on $\mathbb T_\alpha\cap\mathbb T_\beta\subset\mathrm{Sym}^g(\Sigma)$ with
\[
\partial x=\sum_{y}\ \sum_{\substack{\phi\in\pi_2(x,y)\\ \mu(\phi)=1}} \#\big(\cM(\phi)/\R\big)\, U^{n_w(\phi)}V^{n_z(\phi)}\, y .
\]
It carries two Maslov gradings with
$\gr_w(x)-\gr_w(y)=\mu(\phi)-2n_w(\phi)$, $\gr_z(x)-\gr_z(y)=\mu(\phi)-2n_z(\phi)$,
$U$ of bidegree $(-2,0)$, $V$ of bidegree $(0,-2)$, and the Alexander grading
$A=\tfrac12(\gr_w-\gr_z)$, for which $U$, $V$ have degree $-1$, $+1$ and $\partial$ has degree $0$.
\end{definition}

$A$ is well defined because two classes in $\pi_2(x,y)$ differ by $k[\Sigma]+\mathcal P$ with
$\mathcal P$ a periodic domain, $n_w(\Sigma)=n_z(\Sigma)=1$, and, normalising $n_w(\mathcal P)=0$,
$n_z(\mathcal P)=[K]\cdot H(\mathcal P)=0$ since $[K]=0\in H_1(Y;\mathbb Q)$.

\emph{Positivity} is what makes $\partial$ polynomial rather than Laurent: a holomorphic disc meets the
complex hypersurfaces $V_w=\{w\}\times\mathrm{Sym}^{g-1}\Sigma$ and $V_z$ non-negatively, so
$n_w(\phi),n_z(\phi)\ge0$ whenever $\cM(\phi)\ne\emptyset$.

\begin{proposition}[$\partial^2=0$]\label{prop:d2}
For a knot (one $w$, one $z$) $\partial^2=0$ on $\CFK_{\F[U,V]}$.
\end{proposition}
\begin{proof}[Sketch]
Ends of the $1$-dimensional spaces $\cM(\phi)/\R$, $\mu(\phi)=2$, are broken flows and boundary
degenerations.  An $\alpha$-degeneration has domain a component of $\Sigma\setminus\boldsymbol\alpha$; since the
$g$ $\alpha$-curves are linearly independent in $H_1(\Sigma)$ there is exactly one, it contains both
$w$ and $z$, and it contributes $UV\cdot x$ (Ozsv\'ath--Szab\'o: the signed count of degenerations
through a point is $1$).  The same holds for $\beta$, so the two contributions cancel.  For links with
$n$ components the analogous count gives
$\partial^2=\sum_i(U_{w_i}V_{z_i}-U_{w_{i+1}}V_{z_i})$-type curvature, which is why Zemke works with
curved complexes \cite{Zemke-link}.
\end{proof}
This point is worth isolating, because it will recur: \emph{the only codimension-one degenerations
other than breaking carry weight $UV$, and the theory is a complex (rather than a curved one)
because two such contributions cancel.}

\subsection{$\CFK^\infty$ and its subquotients}
Write $[x,i,j]:=U^{-i}V^{-j}x$ for $j-i=A(x)$; these span the $A=0$ part of
$\CFK_{\F[U,V]}\otimes\F[U^{\pm1},V^{\pm1}]$, which is $\CFK^\infty$:
\[
\partial[x,i,j]=\sum \#\big(\cM(\phi)/\R\big)\,[y,\,i-n_w(\phi),\,j-n_z(\phi)],
\]
a $\Z\oplus\Z$-filtered complex of free $\F[U,U^{-1}]$-modules ($U$ acting by $[x,i,j]\mapsto[x,i-1,j-1]$).
Its subquotients recover everything:
\begin{center}\small
\begin{tabular}{ll}
\toprule
$C\{i\le0\}$, $C\{i\ge0\}$, $C\{i=0\}$ & $\CF^-(Y)$, $\CF^+(Y)$, $\widehat{\CF}(Y)$ (basepoint $w$)\\
$C\{j\le0\}$, $C\{j=0\}$ & $\CF^-(Y)$, $\widehat\CF(Y)$ (basepoint $z$)\\
$C\{i=0,j=0\}$, $H_*$ of $C\{i\le 0,\ j=0\}$ & $\widehat{\CFK}$; $\HFK^-$ ($=H_*(\CFK_{\F[U,V]}/V)$)\\
$A^+_s=C\{\max(i,j-s)\ge0\}$, $A^-_s=C\{\min(i,j-s)\le0\}$ & large surgery complexes\\
\bottomrule
\end{tabular}
\end{center}

\begin{proposition}[the two localisations]\label{prop:loc}
$\CFK_{\F[U,V]}\otimes_{\F[V]}\F[V^{\pm1}]\simeq \CF^-(Y,w)\otimes_{\F}\F[V^{\pm1}]$ and
$\CFK_{\F[U,V]}\otimes_{\F[U]}\F[U^{\pm1}]\simeq \CF^-(Y,z)\otimes_{\F}\F[U^{\pm1}]$.
\end{proposition}
\begin{proof}
Rescale generators $x\mapsto V^{A(x)}x$; the differential becomes $\sum U^{n_w}V^{n_z-n_w+A(x)-A(y)}=\sum U^{n_w}$
times a unit-free factor $V^0$, which is the differential of $\CF^-(Y,w)$ tensored with $\F[V^{\pm1}]$.
\end{proof}
So $\CFK$ is a single complex sitting over two copies of $\CF^-(Y)$, one for each basepoint, and the
knot is the \emph{relative position} of these two copies inside $\CF^\infty$.

\subsection{Structure and invariants}
For $K\subset S^3$, $\CFK_{\F[U,V]}$ is a \emph{knot-like complex}: free, finitely generated, with
$H_*(C\otimes\F[U,V^{\pm1}])\cong\F[U,V^{\pm1}]$ supported in $\gr_w=0$ (and symmetrically).
Concordance and surgery invariants are read off from it:
\begin{itemize}[leftmargin=2em]
\item $\tau(K)=-\max\{A(x): x\in \HFK^-(K)\ \text{homogeneous},\ U^nx\neq0\ \forall n\}$;
for $T_{2,3}$ (generators $a,b,c$, $A=1,0,-1$, $\partial b=Ua+Vc$) this gives $\tau=1$.
\item $V_s$ = the $U$-tower drop of $H_*(A^+_s)\to H_*(B^+)$, $\nu^+(K)=\min\{s\ge0:V_s=0\}$, and
for $n>0$, $d(S^3_n(K),i)=d(L(n,1),i)-2\max(V_i,V_{n-i})$ (Ni--Wu \cite{NiWu}).
\item Hom's $\varepsilon$ \cite{Hom}, the local-equivalence group and its homomorphisms $\varphi_j$
(Dai--Hom--Stoffregen--Truong \cite{DHST}), Hendricks--Manolescu's involutive $\iota_K$ \cite{HM-inv} (a skew-filtered map with
$\iota_K^2\simeq\varsigma$, the Sarkar map).
\item The $t$-modified complex (Ozsv\'ath--Stipsicz--Szab\'o \cite{OSS}): over $\F[v^{1/n}]$, $t\in[0,2]\cap\frac1n\Z$,
\[
\partial_t x=\sum \#\big(\cM(\phi)/\R\big)\, v^{\,t\,n_z(\phi)+(2-t)\,n_w(\phi)}\,y,\qquad \gr_t=\gr_w-tA,
\]
and $\Upsilon_K(t)$ = the maximal $\gr_t$ of a $v$-nontorsion class.  $\Upsilon_K(t)=\Upsilon_K(2-t)$,
$\Upsilon_K'(0)=-\tau(K)$, and $\Upsilon_K(t)=(1-|1-t|)\,\sigma(K)/2$ for alternating $K$.
At $t=0$ and $t=2$ one recovers $\CF^-(Y,w)$, $\CF^-(Y,z)$ with $U=v^2$; $t=1$ uses the
$\delta$-type grading $\gr_w-A=\frac12(\gr_w+\gr_z)$.
\end{itemize}
The $t$-family is a one-parameter interpolation between the two basepoints: the exponent is the
pairing of $(n_w,n_z)$ with the covector $(2-t,t)$.  Since $n_w,n_z\ge0$ for all discs, every
$\partial_t$ is defined; conversely, a family of complexes with weights $v^{(2-t)a+tb}$ defined for all
$t\in[0,2]$ forces $(a,b)$ into $\bigcap_t\{(2-t)a+tb\ge0\}=\{a\ge0,b\ge0\}$.  \emph{Positivity on the
quadrant is equivalent to positivity on every ray of the $t$-family.}  This is the form in which
positivity will reappear in gauge theory (\S\ref{sec:proposal}).

\subsection{Surgery}
\emph{Large surgery}: for $n\ge2g(K)-1$ and $|s|\le n/2$, $\CF^\pm(S^3_n(K),[s])\simeq A^\pm_s$.
\emph{Mapping cone}: $\CF^+(S^3_n(K))\simeq\Cone\big(\bigoplus_sA^+_s\xrightarrow{v+h}\bigoplus_sB^+_s\big)$,
with $v_s:A^+_s\to B^+_s=C\{i\ge0\}$ the projection and $h_s:A^+_s\to C\{j\ge s\}\cong C\{j\ge0\}\simeq B^+_{s+n}$,
the last equivalence being induced by the Heegaard moves that carry $z$ to $w$: the \emph{flip map}.
Its existence uses naturality (Juh\'asz--Thurston--Zemke \cite{JTZ}); in the mapping cone it is the only place where
the two basepoints are identified.

\subsection{Functoriality}
Zemke \cite{Zemke-link} assigns $\F[U,V]$-equivariant maps to decorated link cobordisms $(W,F,\mathcal A)$:
$F\subset W$ is an embedded surface and $\mathcal A\subset F$ a collection of arcs dividing $F$ into
$F_w\sqcup F_z$.  With $\tilde\chi(F_w)=\chi(F_w)-\frac12(|\mathbf w_{\rm in}|+|\mathbf w_{\rm out}|)$,
\begin{align*}
\gr_w(F(x))-\gr_w(x)&=\tfrac14\big(c_1(\mathfrak s)^2-2\chi(W)-3\sigma(W)\big)+\tilde\chi(F_w),\\
\gr_z(F(x))-\gr_z(x)&=\tfrac14\big(c_1(\mathfrak s-PD[F])^2-2\chi(W)-3\sigma(W)\big)+\tilde\chi(F_z),\\
A(F(x))-A(x)&=\tfrac12\big(\langle c_1(\mathfrak s),[F]\rangle-[F]\cdot[F]\big)+\tfrac12\big(\tilde\chi(F_w)-\tilde\chi(F_z)\big)
\end{align*}
\cite{Zemke-grading}.  The maps $\Phi=\partial_U\partial$ and $\Psi=\partial_V\partial$ (formal derivatives)
are the maps of the elementary dividing-arc moves, and Sarkar's formula says that moving the basepoints
once around $K$ acts by $\varsigma\simeq\mathrm{id}+\Phi\Psi$ \cite{Sarkar}.

\subsection{Bordered reading}
For $M=Y\setminus\nu K$, Hanselman--Rasmussen--Watson \cite{HRW} encode $\widehat{\CF}(M)$ as an immersed
multicurve $\gamma_M$ in the \emph{marked torus} $\partial M\setminus\{z\}$; Dehn filling is Lagrangian
Floer homology with a line of the filling slope, and $\widehat{\HFK}$ is obtained by pairing with a
meridian placed \emph{at} the marked point so that the two regions adjacent to the marking carry $w$ and
$z$.  For $K\subset S^3$, $\gamma_M$ is determined by $\CFK_{\F[U,V]}/(UV)$; the $UV$-information
(discs covering the marking) is recovered by Hanselman's decorated curves with a bounding chain
\cite{Hanselman}.  The essential component of $\gamma_M$ (homotopic to the longitude, passing between
the lifts of the marking) carries $\tau$ and $\varepsilon$.

%======================================================================
\section{Anatomy: what actually produces $\CFK_{\F[U,V]}$}\label{sec:mech}
%======================================================================

Heegaard Floer homology is itself the output of an Atiyah--Floer heuristic: stretch the
Seiberg--Witten equations on $Y=U_\alpha\cup_\Sigma U_\beta$ along a neck $[-T,T]\times\Sigma$; on
$\R_t\times[-T,T]_s\times\Sigma\cong \C\times\Sigma$ solutions become vortices, the vortex moduli space of
degree $g$ is $\mathrm{Sym}^g\Sigma$, the handlebodies contribute Lagrangians $\mathbb T_\alpha,\mathbb T_\beta$, and
four-dimensional trajectories become holomorphic strips.  We read every ingredient of
\S\ref{sec:cfk} through this lens, because a gauge-theoretic $\CFK^\infty$ should be designed so that
its adiabatic limit \emph{is} a complex of the shape of \S\ref{sec:cfk}.

\subsection{Mechanisms}
\begin{enumerate}[label=(M\arabic*),leftmargin=2.6em]
\item\label{M:charges} \textbf{Two non-negative charges.}  $n_w(\phi),n_z(\phi)\in\Z_{\ge0}$, from positivity
of intersections with complex divisors $V_w,V_z$ disjoint from $\mathbb T_\alpha\cup\mathbb T_\beta$.
In vortex language, $V_w$ is the set of vortex configurations with a vortex at $w$; for a
four-dimensional trajectory in the neck, $n_w$ counts zeros of the \emph{vortex component} $\alpha$
of the spinor on the strip $\R_t\times[-T,T]_s\times\{w\}$, which is a complex curve in $\C\times\Sigma$.
\item\label{M:alex} \textbf{One of the charges is topological.} $n_w-n_z$ depends only on endpoints
($=A(x)-A(y)$), because $V_w-V_z$ is null-homologous in $\mathrm{Sym}^g\Sigma$ relative to the Lagrangians
(\S\ref{sec:cfk}).  The other combination $n_w+n_z$ is energy-like.
\item\label{M:quadrant} \textbf{Positivity on a quadrant.} Equivalently, positivity of every member of
the $t$-family of energies $(2-t)n_w+tn_z$, $t\in[0,2]$.
\item\label{M:degen} \textbf{Codimension-one degenerations carry $UV$ and cancel} (Proposition~\ref{prop:d2}).
\item\label{M:loc} \textbf{Two localisations}, each recovering the closed $3$-manifold invariant
(Proposition~\ref{prop:loc}).  Forgetting $z$ allows Heegaard moves across $z$, which is why
$V=1$ gives an invariant of $Y$ alone.
\item\label{M:flip} \textbf{A flip map} identifying the two localisations, and the Sarkar map
$\varsigma\simeq 1+\Phi\Psi$ as its square (monodromy of the basepoints once around $K$).
\item\label{M:cob} \textbf{Cobordisms with dividing sets}, with grading shifts additive in
$\tilde\chi(F_w)$, $\tilde\chi(F_z)$.
\item\label{M:surg} \textbf{Large surgery and the mapping cone}, which use \ref{M:loc} and \ref{M:flip}.
\end{enumerate}

\subsection{The geometric position of the knot}
In the neck picture the knot meets $\Sigma$ transversally at $w$ (upward) and $z$ (downward), and
$K=\gamma_w\cup\overline{\gamma_z}$ with $\gamma_w,\gamma_z$ the upward flow arcs from $p_0$ to $p_3$.
On $[-T,T]\times\Sigma$ the vortex splitting $S=E\oplus EK_\Sigma^{-1}$ is the $\pm i$ eigenspace splitting
of Clifford multiplication by the neck direction $\partial_s$.  Along $K$ one has the intrinsic splitting
\[
S|_K=E_+(\dot K)\oplus E_-(\dot K),\qquad E_\pm(\dot K)=\ker\big(\rho(\dot K)\mp i\big),
\]
and since $\dot K=\partial_s$ at $w$ and $\dot K=-\partial_s$ at $z$, \emph{the vortex component is the
$E_+(\dot K)$ part of $\Phi$ near $w$ and the $E_-(\dot K)$ part near $z$}.  So $n_w$ and $n_z$ count zeros of
\emph{opposite} chiral components of the spinor along the two arcs of $K$, with the switch occurring at
$p_0,p_3$, where the two arcs meet.  This is the first concrete clue to what the gauge-theoretic charges
must be (\S\ref{subsec:monopole}), and it explains why $U$ and $V$ are exchanged by reversing $K$.

\subsection{Dictionary}
\begin{center}
\small
\begin{tabular}{p{4.3cm}p{4.5cm}p{4.9cm}}
\toprule
Heegaard Floer & Atiyah--Floer intermediary & Gauge theory (proposed)\\
\midrule
$\mathrm{Sym}^g\Sigma$; $\mathbb T_\alpha,\mathbb T_\beta$ & moduli of the Heegaard surface; Lagrangians of handlebody boundary values & configuration space of $Y$ with a neck of length $T$\\
basepoints $w,z\in\Sigma$ & two punctures with meridional holonomy & codimension-two singular locus $K$ (holonomy parameter)\\
$V_w,V_z$ & two divisors/nef classes $\xi_w,\xi_z$ & two topological charges (\S\ref{sec:proposal})\\
$n_w-n_z$ topological & $\xi_w-\xi_z$ relatively trivial (monopole); flag classes (instanton) & flux of curvature through $\R\times K$; monopole number $l$\\
$t$-family, $t\in[0,2]$ & Kähler classes $[\omega_\alpha]$, $\alpha\in(0,\tfrac12)$ & holonomy parameter; $t=4\alpha$\\
boundary degenerations ($UV$) & Maslov-$2$ discs on $L_\alpha,L_\beta$ & bubbles at $\R\times K$ and instanton bubbles\\
moves avoiding $w,z$ & Lagrangian correspondences of handle attachments & metric/perturbation changes fixed near $K$\\
dividing arcs on $F$ & changes of puncture type & creases of the edge boundary condition\\
\bottomrule
\end{tabular}
\end{center}

%======================================================================
\section{What exists on the gauge side, and the gap}\label{sec:landscape}
%======================================================================

\begin{itemize}[leftmargin=2em]
\item \textbf{Hat level.} $KHM$, $KHI$ \cite{KM-sutures}; $\chi(KHI)=\Delta_K$ \cite{KM-alex,Lim};
$KHM\cong\widehat{\HFK}$ (Lekili together with $\HF=\HM$ \cite{Lekili,KLT}; Baldwin--Sivek \cite{BS-equiv};
Colin--Ghiggini--Honda directly \cite{CGH-sutured}).  $I^\natural(K)\cong KHI(K)$ \cite{KM-unknot},
where $I^\natural$ is singular instanton homology with holonomy $\mathrm{diag}(i,-i)$, i.e.\ parameter
$\tfrac14$ (the monotone value \cite{KM-knothom}).
\item \textbf{One-variable versions.} Li's $KHM^-$, $KHI^-$ (direct limits of sutured homologies
\cite{Li-minus}); Ghosh--Zemke's free model over $\C[U]$ as a cone of sutured complexes \cite{GZ};
Daemi--Scaduto's $S$-complex $\tilde C(Y,K)$ from framed singular connections, whose $\hat C$
is a counterpart of $\HFK^-$, together with a universal local system over $\Z[U^{\pm1},T^{\pm1}]$, $U$
recording Chern--Simons energy and $T$ holonomy (monopole charge) \cite{DS-equiv}; filtered
versions \cite{DS-clasp,DISST}; Miller Eismeier's $I^{\pm,\infty}$ \cite{ME,DME}.
\item \textbf{Surgery.} $\nu^\sharp,\tau^\sharp$ and the dimension formula for $I^\sharp(S^3_{p/q}(K))$
\cite{BS-conc}; Li--Ye large and integral surgery formulae \cite{LY-large,LY-surgery}, built on
differentials $d_i^j$ between Alexander gradings of $KHI$ and ``bent complexes'' $A_s$.
\item \textbf{Bordered/Atiyah--Floer.} Pillowcase homology for 2-tangles
\cite{HHK1,HHK2,CHKK}; Donghao Wang's monopole Floer homology for toroidal boundary \cite{Wang}
(recovering monopole link Floer homology at the hat level); Atiyah--Floer for $SO(3)$ \cite{DFL};
the instanton theory with Lagrangian boundary conditions \cite{SW-lag}.
\end{itemize}

\textbf{The gap.}  No gauge theory produces a complex over a two-variable polynomial ring, a
$\Z\oplus\Z$-filtered $\CFK^\infty$, or a flip map.  Li--Ye list the differential $D$ on
$\CFK^\infty$ (with $D^2=0$) among the Heegaard Floer structures and state that they do not know how to
construct its analogue \cite{LY-large}.  Daemi--Scaduto's two-variable \emph{Laurent} local system is the
nearest structure: it parametrises the right lattice, but carries no positivity and hence no
quadrant/filtration structure.  The proposal below is designed to fill exactly this gap: produce the
quadrant (mechanisms \ref{M:charges}--\ref{M:quadrant}) from an Atiyah--Floer degeneration.

%======================================================================
\section{The proposal: an Atiyah--Floer construction of the quadrant}\label{sec:proposal}
%======================================================================

\subsection{Principle}
The construction has three layers.
\begin{enumerate}[label=(\Roman*),leftmargin=2.4em]
\item A \textbf{symplectic target}: the moduli space $\mathcal M$ of the Heegaard surface punctured at
$K\cap\Sigma=\{w,z\}$, Lagrangians $L_\alpha,L_\beta$ from the two handlebody--arc pairs, and two
\emph{nef} classes $\xi_w,\xi_z$ whose positive combinations are the symplectic classes.  This is the
analogue of $(\mathrm{Sym}^g\Sigma,\mathbb T_\alpha,\mathbb T_\beta,V_w,V_z)$.
\item A \textbf{gauge complex} on $\R\times(Y,K)$ whose trajectories are weighted by
$U^{c_w(A)}V^{c_z(A)}$ for two \emph{topological charges} $c_w,c_z$ that are defined for every
configuration, without reference to a splitting.
\item An \textbf{adiabatic comparison}: for a metric with neck length $T$ along $\Sigma$ and energy
$\le E$, trajectories converge to holomorphic strips in $\mathcal M$ and $c_w,c_z$ converge to
$\xi_w,\xi_z$ evaluated on the limit.  Since $c_w,c_z$ are integers, approximate positivity becomes exact
positivity:
\end{enumerate}

\begin{lemma}[integrality bootstrap]\label{lem:bootstrap}
Let $c$ be a $\Z$-valued function on trajectories that is locally constant on moduli spaces.  Suppose that
for every $E$ there is $T(E)$ such that for $T\ge T(E)$ every trajectory $A$ of energy $\le E$ satisfies
$c(A)\ge \xi(u_A)-\epsilon(T)$, where $u_A$ is a (possibly broken, bubbled) holomorphic limit, $\xi(u_A)\ge0$
and $\epsilon(T)\to0$.  Then $c(A)\ge0$ for all such $A$ once $\epsilon(T)<1$.
\end{lemma}
The proof is one line, but the lemma fixes the architecture.  The complex is defined at finite $T$, with
an energy truncation, and the full complex is the direct limit over $E\to\infty$ with $T\ge T(E)$,
as in Taubes' passage from $\HM$ to ECH through action-filtered complexes.  All the analysis is in the
hypothesis: uniform convergence of the charges in the adiabatic limit, including over bubbles.

%----------------------------------------------------------------------
\subsection{What a gauge theory must supply}\label{subsec:requirements}
%----------------------------------------------------------------------
\begin{lemma}[$\CFK$ is $\CF^-(Y)$ plus one function and one inequality]\label{lem:reform}
Let $\mathcal H$ be doubly pointed, and let $\CF^-(\mathcal H,w)$ be the free $\F[U]$-module on generators with
$\partial x=\sum_\phi\#\big(\cM(\phi)/\R\big)\,U^{n_w(\phi)}y$.  Then $\CFK_{\F[U,V]}(\mathcal H)$ is obtained by
replacing each $U^{n_w(\phi)}$ with $U^{n_w(\phi)}V^{\,n_w(\phi)-A(x)+A(y)}$.  The knot is thus encoded by:
\begin{enumerate}[label=(\roman*)]
\item a function $A$ from generators to $\Z$, and
\item the inequality $n_w(\phi)\ge A(x)-A(y)$ whenever $\cM(\phi)\ne\emptyset$, which is $n_z(\phi)\ge0$.
\end{enumerate}
Moreover $\CFK_{\F[U,V]}\otimes\F[U,V]/(U-1,V-1)\cong\CF^-(\mathcal H,w)\otimes\F[U]/(U-1)$, so this specialisation does
not depend on $K$.
\end{lemma}
\begin{proof}
$n_z=n_w-(A(x)-A(y))$ (\S\ref{sec:cfk}); the last claim is immediate from the substitution.
\end{proof}

Read as a specification, Lemma~\ref{lem:reform} gives three requirements.
\begin{enumerate}[label=(Q\arabic*),leftmargin=2.6em]
\item\label{Q:base} The generators and moduli spaces are those of a Floer theory of \emph{the closed manifold} $Y$,
for auxiliary data adapted to $K$.  The knot enters only through weights.  This excludes the singular theories
(they change the generators), and is why \S\ref{subsec:symp} is only a variant.
\item\label{Q:alex} An integer Alexander function on generators.
\item\label{Q:pos} One-sided positivity of a $z$-charge $n_z$ whose difference from the $w$-charge is
$A(x)-A(y)$.
\end{enumerate}
In monopole and instanton theory, $U$ comes from equivariance at the reducible, not from a weight on
trajectories.  The natural gauge-theoretic form is therefore the \emph{dual} one: keep $\CF^-(Y)$ as Floer theory
already produces it, an $S^1$- or $SO(3)$-equivariant complex whose $U$ is the point class at $w$, and add a single
variable $V$ weighting trajectories by a non-negative $z$-charge.  The Heegaard Floer statement corresponding to
this is literally Lemma~\ref{lem:reform}.

%----------------------------------------------------------------------
\subsection{The monopole model}\label{subsec:monopole}
%----------------------------------------------------------------------
\paragraph{Data.}
Take a doubly pointed Heegaard splitting with $K=\gamma_w\cup\overline{\gamma_z}$ as in \S\ref{sec:mech}, and a metric
$g_T$ that is a product on a neck $N_T=[-T,T]_s\times\Sigma$.  Then $K\cap N_T=[-T,T]\times\{w,z\}$, with $\dot K=\partial_s$ at $w$
and $\dot K=-\partial_s$ at $z$.  Use Kronheimer--Mrowka perturbations \cite{KM-book} supported in the handlebodies, plus the standard
neck perturbation (an area form on $\Sigma$, playing the role of the vortex parameter).  Write
$\gamma_z^T=\{z\}\times[-T,T]$ and $\gamma_w^T=\{w\}\times[-T,T]$.

\paragraph{Charges.}
For a configuration $\gamma=(A,s,\phi)$ on $\R\times Y$ in the blown-up model ($\|\phi(t)\|_{L^2}=1$), let
\[
\phi^-_z:=\pi_{E_-(\dot K)}\big(\phi|_{\R\times\gamma_z^T}\big),\qquad
\phi^+_w:=\pi_{E_+(\dot K)}\big(\phi|_{\R\times\gamma_w^T}\big),
\]
the vortex components of \S\ref{sec:mech}.  They are sections of complex line bundles over the strips
$\R\times\gamma^T_{z}$ and $\R\times\gamma^T_{w}$; trivialise these bundles by pulling back from the arcs.

\begin{definition}[charges]
Say $\gamma$ is \emph{admissible} if $\phi^\pm$ do not vanish on $\R\times\{\pm T\}$, nor at the limiting critical points along
the arcs.  For admissible $\gamma$, set $n_z(\gamma):=\frac1{2\pi}\Delta\arg\phi^-_z$ around $\partial\big([-R,R]\times\gamma^T_z\big)$
as $R\to\infty$, and define $n_w$ likewise.  Equivalently, these are the signed zero counts on the strips.
\end{definition}
On admissible moduli spaces $n_z$ is locally constant and additive under breaking and gluing: the contributions of the
ends $t=\pm R$ cancel in pairs.  Admissibility is the gauge-theoretic shadow of $V_z\cap(\mathbb T_\alpha\cup\mathbb T_\beta)=\emptyset$.
At $s=\pm T$ the configuration is close to a handlebody boundary value whose vortex divisor lies on the $\alpha$- or
$\beta$-curves, and $z$ does not.

\begin{definition}[monopole knot complex; conditional on Conjecture~\ref{conj:Mpos}]\label{def:CFKHM}
$\CFK^{\HM}_T(Y,K):=\widehat{CM}(Y;g_T,\mathfrak q)\otimes_\F\F[V]$.  Each matrix entry of the Kronheimer--Mrowka
differential, including the boundary-obstructed components $\bar\partial^s_u$ and so on, is replaced by its count weighted
by $V^{n_z}$.  The checked and barred versions are defined similarly.
\end{definition}

\begin{conjecture}[admissibility and positivity in the adiabatic regime]\label{conj:Mpos}
For every $E$ there is $T(E)$ such that, for $T\ge T(E)$ and generic $\mathfrak q$, every trajectory of energy $\le E$ in the
moduli spaces of dimension $\le1$ is admissible, and satisfies $n_z\ge0$, $n_w\ge0$ and
$n_w-n_z=A(a)-A(b)$ for an integer function $A$ on critical points.
\end{conjecture}
Strategy: the adiabatic limit $T\to\infty$ (\S\ref{sec:analysis}) presents a trajectory as a vortex strip in
$\mathrm{Sym}^g\Sigma$ glued to handlebody solutions.  On the neck, $\phi^-_z$ converges, after rescaling, to the
holomorphic section whose zeros are the intersections with $V_z$.  Positivity of zeros of a nearly holomorphic
section (Lemma~\ref{lem:posvortex}) and the integrality bootstrap (Lemma~\ref{lem:bootstrap}) then give the
inequalities, and avoidance at $s=\pm T$ gives admissibility.

\begin{theorem}[structure, assuming Conjecture~\ref{conj:Mpos}]\label{thm:Mstructure}
\begin{enumerate}[label=(\alph*)]
\item $\partial_V^2=0$.  The ends of $1$-dimensional moduli spaces, including boundary-obstructed ones, are exactly
those used by Kronheimer--Mrowka, and $n_z$ is additive over them.  No new ends appear, because admissibility is
uniform.
\item $\CFK^{\HM}_T\otimes\F[V]/(V-1)=\widehat{CM}(Y)$, so $V=1$ recovers $\widehat{\HM}(Y)\cong\HF^-(Y)$ (localisation \ref{M:loc}
is built in).
\item The $U$-map, defined by the codimension-two condition $\phi^+_w(0,w)=0$, i.e.\ a vortex at $w$, commutes with
$\partial_V$ up to $V$-weighted homotopy.  So $H_*(\CFK^{\HM}_T)$ is an $\F[U,V]$-module.
\end{enumerate}
\end{theorem}
\begin{proof}[Proof sketch]
(a) Because $n_z$ is a locally constant integer on admissible moduli spaces, and additive under gluing, each end of a
compactified $1$-dimensional space contributes $V^{n_z}$ times the corresponding product.  Kronheimer--Mrowka's
identities among the components then hold verbatim over $\F[V]$.  Positivity makes all exponents non-negative.
(b) is immediate.  (c) The usual proof that $U$ is a chain map uses $1$-dimensional moduli cut down by a codimension-$2$
condition; weighting by $V^{n_z}$ commutes with that argument.
\end{proof}

\begin{conjecture}[knot Atiyah--Floer, monopole form]\label{conj:MAF}
For $T\ge T(E)$, and after the direct limit over $E$, $\CFK^{\HM}_T(Y,K)\simeq\CFK_{\F[U,V]}(\mathcal H)$ as complexes over $\F[V]$,
with $U$ corresponding to the point-class map at $w$.  The homotopy type is independent of the doubly pointed splitting,
$T$ and $\mathfrak q$.  In particular it is an invariant of $(Y,K)$.
\end{conjecture}
Independence of $\mathfrak q$ and $T$ comes from continuation maps that are $V$-weighted.  These are well defined as soon
as admissibility holds along $1$-parameter families, which is again uniform positivity.  Independence of the splitting
is the gauge-theoretic counterpart of invariance under moves avoiding $w,z$.  Stabilisation and handleslides away from
$K$ are handled by neck-stretching along the extra surfaces (\S\ref{subsec:inv}).

\paragraph{The $U$-localisation is a real prediction.}
In Heegaard Floer, inverting $U$ gives $\CF^\infty$ with the $j$-filtration, whose $j\le0$ part is $\CF^-(Y,z)$.
In the monopole model, inverting $U$ passes to $\overline{CM}(Y)$, which is generated by the reducible's eigenvectors.
The prediction is that $\overline{CM}(Y)\otimes\F[V]$ with $V$-weights recovers $\widehat{\HM}(Y)$ with $V$ in the role of $U$.  The
$V$-weight of a reducible flow from $\psi_\lambda$ to $\psi_\mu$ is the difference of windings of the $E_-(\dot K)$-components of the
Dirac eigenvectors along $\gamma_z$.  So the conjecture predicts that \emph{edge windings of Dirac eigenvectors along an arc
of $K$ encode $\HM(Y)$ seen from $z$}.  For $(S^3,\text{unknot})$ with the round metric and $K$ a great circle, the
$T^2$-symmetry of $S^3\subset\C^2$ makes this an explicit separation-of-variables computation (\S\ref{sec:tests}).

%----------------------------------------------------------------------
\subsection{The instanton model}\label{subsec:instanton}
%----------------------------------------------------------------------
\paragraph{Symplectic side.}
Atiyah--Floer replaces $\mathrm{Sym}^g\Sigma$ with $\mathcal R(\Sigma)$, the $SU(2)$ moduli space of the \emph{unpunctured} Heegaard surface.
Requirement \ref{Q:base} forbids punctures.  Framing at a point in a handlebody, or extended moduli spaces as in
Manolescu--Woodward \cite{MW}, removes the reducible singularities.  The handlebody Lagrangians are $L_\alpha,L_\beta$.  The
analogue of $V_p$ is a \emph{theta divisor attached to the point $p$}.  Under the Abel--Jacobi map,
$V_p=\{D\ge p\}\subset\mathrm{Sym}^g\Sigma$ is the translate of $W_{g-1}$ by $p$, and the rank-two analogue is the generalised
theta divisor
\[
D_p:=\Theta_{F_0(p)}=\{E:\ H^0(E\otimes F_0(p))\ne0\},\qquad \deg F_0=g-2,
\]
on the moduli of semistable bundles, transported to $\mathcal R(\Sigma)$ by Narasimhan--Seshadri.  Then $[D_w]=[D_z]=\theta$, the
positive generator of $H^2(\mathcal R(\Sigma))$, as $[V_w]=[V_z]$ in Heegaard Floer, and $D_w-D_z$ is the twist by
$\mathcal O(w-z)$: a unit $U(1)$-flux through $\Sigma$ at $w$ and $-1$ at $z$.  That is exactly the flux of a line bundle
whose curvature is Poincaré dual to $K$.  Since $c_1(T\mathcal R(\Sigma))=2\theta$, a disc with $u\cdot D_w=1$ changes the grading by $4$
in instanton units, which is $2$ in Heegaard Floer units.  This matches $\deg U=4$ for the $SO(3)$-equivariant
theory \cite{ME,DME} under the usual doubling of instanton gradings.

\paragraph{Where the knot enters.}
There is one point where the instanton picture differs essentially from $\mathrm{Sym}^g\Sigma$.  In $\mathcal R(\Sigma)$ the
Lagrangian $L_\alpha$ is canonical: it depends on the handlebody, not on the $\alpha$-curves.  In Heegaard Floer, by contrast,
$\mathbb T_\alpha$ depends on the curves up to isotopy \emph{in the complement of $V_w,V_z$}, and that dependence carries the
knot (moving an $\alpha$-curve across $w$ changes $K$).  So the arcs $a_\alpha,a_\beta$ cannot enter through $L_\alpha,L_\beta$.  They
enter through the \emph{trivialisations} of the determinant line $\mathcal D_p=\det\bar\partial_{\,\cdot\,\otimes F_0(p)}$ over $L_\alpha$ and $L_\beta$.
Along $L_\alpha$ the line $\mathcal D_z$ is trivialised by the three-dimensional Dirac operator on $U_\alpha$ coupled to the flat
connection and to a $U(1)$-connection whose curvature is the Thom form of the part of $\gamma_z$ inside $U_\alpha$.  That
connection has a Dirac-monopole source at $p_0$.  The charge of a strip $u$ is the relative degree
\[
c_z(u):=\big\langle c_1(\mathcal D_z;\ \tau_\alpha,\tau_\beta),\,[u]\big\rangle\in\Z,
\]
which is topological and knot-dependent.  $D_z$ enters only as a device for positivity.

\begin{problem}[compatible divisors]\label{prob:avoid}
Find holomorphic sections of $\mathcal D_z^{\otimes k}$ (and $\mathcal D_w^{\otimes k}$) over the extended $\mathcal R(\Sigma)$ that do not vanish on
$L_\alpha\cup L_\beta$ and are homotopic there, through non-vanishing sections, to $\tau_\alpha^{\otimes k}$ and $\tau_\beta^{\otimes k}$.
Their zero sets $D_z$ then satisfy $c_z(u)=\tfrac1k\,u\cdot D_z\ge0$ for holomorphic strips.  (A necessary condition, $\theta|_{L}=0$,
holds because handlebody Lagrangians are Bohr--Sommerfeld for the Chern--Simons bundle.)  Failing an explicit
construction, use Donaldson divisors in the complement of $L_\alpha\cup L_\beta$ in the prescribed relative class, following \cite{AGM,CW}.
\end{problem}

\paragraph{Gauge side: the $z$-charge as a determinant-line crossing number.}
For an $SU(2)$ connection $\mathbb A$ on $\R\times Y$ with neck $N_T$, the slices $\mathbb A|_{\{(t,s)\}\times\Sigma}$ define a map from the strip
$\R_t\times[-T,T]_s$ to the space of connections on $\Sigma$, and the twisted Dolbeault operators
$\bar\partial_{\mathbb A(t,s)\otimes F_0(z)}$ form a family of index zero.  Their Quillen determinant line carries the canonical
section $\det\bar\partial$, whose zero set over $\mathcal R(\Sigma)$ is $D_z$.  Define $c_z(\mathbb A)$ as the relative degree of the pulled-back line $\mathcal D_z$ over the strip, with respect to the handlebody
trivialisations $\tau_\alpha,\tau_\beta$ at $s=\mp T$ and the critical points at $t=\pm\infty$.  When a compatible section exists, it is
the number of $(t,s)$ with $\ker\bar\partial_{\mathbb A(t,s)\otimes F_0(z)}\ne0$, counted with orientation: a \emph{spectral crossing count},
the instanton analogue of $n_z$.
By the families index theorem in the adiabatic limit it equals a relative index of the four-dimensional Dirac operator
coupled to $\mathbb A\otimes\lambda_z$, where $\lambda_z$ is a $U(1)$-connection with curvature the Thom form of $\R\times\gamma_z$.  This
connection has Dirac-monopole sources on the lines $\R\times\{p_0\}$ and $\R\times\{p_3\}$, where the two arcs of $K$ meet.
The difference $c_w-c_z$ then corresponds to the source-free connection with curvature Thom$(\R\times K)$, so it is an
endpoint quantity.  This is \ref{Q:alex}: an \emph{Alexander function given by a relative spectral flow of the Dirac
operator twisted by unit flux along $K$}.

\begin{definition}[instanton knot complex; conditional on positivity, e.g.\ via Problem~\ref{prob:avoid}]
$\CFK^{I}_T(Y,K):=$ the $SO(3)$-equivariant instanton complex of $Y$ (Miller Eismeier; Daemi--Miller Eismeier), with its
equivariant variable $U$, for the metric $g_T$, tensored with $\F[V]$ and with trajectories weighted by $V^{c_z}$.  Its
symplectic counterpart is $CF(L_\alpha,L_\beta)$ in the extended $\mathcal R(\Sigma)$ with weights $U^{u\cdot D_w}V^{u\cdot D_z}$.
\end{definition}
The analogues of Conjectures~\ref{conj:Mpos} and~\ref{conj:MAF} are stated in the same way.  Two features are specific
to instantons and are worth flagging:
\begin{itemize}[leftmargin=2em]
\item \emph{Deck versus equivariant variable.}  In Heegaard Floer, $U$ is both the deck transformation by $[\Sigma]$ and the
point-class map.  In instanton theory the deck generator (degree $8$) and the equivariant class (degree $4$) are
different, so the analogue of ``the $U$-map equals multiplication by $U$'' is a theorem to prove.  The prediction is that
$U_{\rm eq}^2$ is homotopic to the deck transformation on the equivariant complex.
\item \emph{Bubbling.}  Uhlenbeck bubbles have $c_w=c_z=$ instanton number.  They are $UV$-type and occur in codimension
$\ge2$ in the monotone setting, which matches the sphere bubbles $k[\Sigma]$ in Heegaard Floer.
\end{itemize}


\subsection{A singular variant: parabolic moduli and the nef segment}\label{subsec:symp}
This subsection and the next develop a \emph{different} two-variable theory, built on singular
instantons.  It is not the analogue of $\CFK$.  Its specialisation $U=V=1$ is $I^\natural(K)$, which
depends on the knot, whereas $\CFK_{\F[U,V]}\otimes\F[U,V]/(U-1,V-1)$ does not (\S\ref{subsec:requirements}).
It is a bivariant refinement of the Kronheimer--Mrowka and Daemi--Scaduto local systems, on the Khovanov
side of the dictionary.  It is included because its positivity mechanism, nefness of the two ends of a
segment of Kähler classes, is the cleanest instance of \ref{M:quadrant}, and because it yields an
instanton analogue of the $t$-family.
Fix a doubly pointed Heegaard splitting of genus $g\ge1$: $K\cap\Sigma=\{w,z\}$ and $K\cap U_\alpha=a_\alpha$,
$K\cap U_\beta=a_\beta$ are boundary-parallel arcs.  For $\alpha\in(0,\tfrac12)$ let
\[
\mathcal M_\alpha=\big\{\text{flat }SU(2)\text{ connections on }\Sigma\setminus\{w,z\},\ \mathrm{hol}_w,\mathrm{hol}_z\sim e^{2\pi i\,\mathrm{diag}(\alpha,-\alpha)}\big\}/\cG,
\]
of dimension $6g-6+2\cdot2=6g-2$, with its Goldman form $\omega_\alpha$.  Let
$L_\alpha^{(\alpha)}\subset\mathcal M_\alpha$ be the connections extending flatly over $U_\alpha\setminus a_\alpha$ with
meridional holonomy in the same class.  Since $U_\alpha\setminus a_\alpha$ is a genus-$(g+1)$ handlebody, its
representation space is $SU(2)^g\times C_\alpha/SU(2)$, of dimension $3g+2-3=3g-1=\tfrac12\dim\mathcal M_\alpha$:
it is Lagrangian.  Define $L_\beta^{(\alpha)}$ likewise.  At $\alpha=\tfrac14$ this is the traceless
setting, which is monotone \cite{KM-knothom}; Horton's symplectic instanton homology of knots works with
traceless character varieties in this spirit \cite{Horton}.

By Mehta--Seshadri \cite{MS}, $\mathcal M_\alpha$ is the moduli space of rank-$2$, trivial-determinant
parabolic bundles on $(\Sigma,\{w,z\})$ with weights $\pm\alpha$ at both points that are stable for the
weight $\alpha$.

\begin{proposition}[no walls for two punctures]\label{prop:nowalls}
For $\alpha\in(0,\tfrac12)$ the $\alpha$-(semi)stability of a parabolic bundle with two marked points is
independent of $\alpha$.  Hence the stable loci $\mathcal M^s_\alpha$ are a single complex manifold $\mathcal M^s$,
and the $\omega_\alpha$ are Kähler forms on it for one fixed complex structure $J$.
\end{proposition}
\begin{proof}
For a line subbundle $L\subset E$ of degree $d$, $\mathrm{pardeg}(L)=d+\alpha(\varepsilon_w+\varepsilon_z)$ with
$\varepsilon_p=+1$ if $L_p$ is the flag at $p$ and $-1$ otherwise.  Stability is $\mathrm{pardeg}(L)<0$ for all $L$.
Since $|\alpha(\varepsilon_w+\varepsilon_z)|\le2\alpha<1$, $\mathrm{pardeg}(L)<0$ iff $d<0$, or $d=0$ and
$\varepsilon_w+\varepsilon_z=-2$; semistability adds $d=0,\ \varepsilon_w+\varepsilon_z=0$.  Neither condition involves $\alpha$.
\end{proof}
\begin{remark}\label{rem:walls}
For a bridge sphere with $2n\ge4$ punctures this fails.  With $n=2$, a degree $-1$ subbundle through all four
flags has $\mathrm{pardeg}=-1+4\alpha$, so $\alpha=\tfrac14$ is itself a wall.  This is the pillowcase, whose
corners are the reducibles \cite{HHK1}.  It is consistent with Heegaard Floer, where multi-pointed diagrams
need one pair of variables per pair of basepoints.  In the present proposal the doubly pointed genus-$g$
splitting is the correct Atiyah--Floer setting, \emph{because} it is the one without walls.
\end{remark}

Through the parabolic theta bundle \cite{Pauly}, the classes $[\omega_\alpha]$ on $\mathcal M^s$ are affine in $\alpha$
(equivalently, Duistermaat--Heckman for the residual torus at the punctures):
\[
[\omega_\alpha]=(1-2\alpha)\,\xi_w+2\alpha\,\xi_z,\qquad \xi_w=\lim_{\alpha\to0}[\omega_\alpha],\quad \xi_z=\lim_{\alpha\to1/2}[\omega_\alpha].
\]
Here $\xi_w$ is pulled back from the moduli space of the underlying bundle (forget the flags; Mehta--Seshadri
at $\alpha\to0$), and $\xi_z$ is pulled back along the double Hecke modification at $w,z$ (at $\alpha\to\tfrac12$
the meridional holonomy tends to $-1$).  Then $\xi_z-\xi_w$ is a combination of the flag classes
$e_w=c_1(\mathcal L_w)$, $e_z=c_1(\mathcal L_z)$, where $\mathcal L_p$ is the eigenline bundle of $\mathrm{hol}_p$.  With
$t:=4\alpha\in(0,2)$,
\[
[\omega_\alpha]=\tfrac12\big((2-t)\,\xi_w+t\,\xi_z\big),
\]
which is the exponent of the Ozsv\'ath--Stipsicz--Szab\'o $t$-modified differential with $(n_w,n_z)$
replaced by $(\xi_w,\xi_z)$.  The monotone traceless value $\alpha=\tfrac14$ is $t=1$, the $\delta$-graded
point of the $t$-family.

\begin{proposition}[nef boundary classes]\label{prop:nef}
For every non-constant $J$-holomorphic sphere $u$ in $\mathcal M^s$, $\xi_w(u)\ge0$ and $\xi_z(u)\ge0$.
\end{proposition}
\begin{proof}
$u$ is holomorphic for $J$, which is compatible with every $\omega_\alpha$ (Proposition~\ref{prop:nowalls}), so
$0<\omega_\alpha(u)=(1-2\alpha)\xi_w(u)+2\alpha\xi_z(u)$ for all $\alpha\in(0,\tfrac12)$.  Let $\alpha\to0$ and $\alpha\to\tfrac12$.
\end{proof}

This is the singular incarnation of \ref{M:quadrant}: \emph{the two variables are the two ends of the segment of
Kähler classes swept out by the holonomy parameter, and positivity of $n_w,n_z$ is nefness of its endpoints.}
Positivity is automatic in Heegaard Floer because $V_w,V_z$ are complex divisors.  Here it is a theorem about
the Kähler cone, and it holds only because there are no walls (Remark~\ref{rem:walls}).

\paragraph{A genuine difference from Heegaard Floer.}
In $\mathrm{Sym}^g\Sigma$, $[V_w]=[V_z]$, so $n_w-n_z$ is topological \ref{M:alex}.  In $\mathcal M^s$, $\xi_z-\xi_w\neq0$:
for a stable underlying bundle $E$, the sphere $P_w=\mathbb P(E_w)$ obtained by moving the flag at $w$ has
$\xi_w(P_w)=0<\xi_z(P_w)$.  Symmetrically, a flag sphere $Q$ on the Hecke side has $\xi_z(Q)=0<\xi_w(Q)$.  By
monotonicity $c_1$ is a positive multiple of $[\omega_{1/4}]=\frac12(\xi_w+\xi_z)$ on spheres.  So if
$m\,\xi_z(P_w)=n\,\xi_w(Q)=N$, the class $m[P_w]-n[Q]$ has $c_1=0$, zero $\omega_{1/4}$-energy, and weight
$U^{-N}V^{N}$.  The instanton complex is therefore homogeneous for the $\delta$-type grading
$\tfrac12(\gr_w+\gr_z)$ but \emph{not} for the Alexander grading, and the index-zero, energy-zero deck classes
are generated by a power of $T:=V/U$.  This matches the extra holonomy variable of Daemi--Scaduto's local system and the rank-two lattice
$(k,l)$ of instanton and monopole numbers (\S\ref{subsec:gauge-inst}).  Alexander homogeneity has to be recovered
separately, for instance from the eigenspace decomposition of $\mu(\hat S)$ as for $KHI$ \cite{KM-sutures}.

\begin{problem}[relative positivity]\label{prob:relative}
Construct closed $2$-forms, or complex divisors $D_w,D_z$, representing $\xi_w,\xi_z$ in
$H^2(\mathcal M,L_\alpha\cup L_\beta)$ (with $L_\bullet=L_\bullet^{(1/4)}$), disjoint from $L_\alpha\cup L_\beta$ and positive on
$J$-holomorphic strips.
\end{problem}
Here is a route.  Handlebody Lagrangians are Bohr--Sommerfeld for the Chern--Simons prequantum bundle: the
Chern--Simons invariant of the flat extension is a covariantly constant section over $L$.  Show the same for the line bundles of $\xi_w$ (extension
over $U_\alpha$ with the arc erased) and $\xi_z$ (extension with holonomy $-1$ around the arc).  Then take
Donaldson hypersurfaces in the complement of the Lagrangians, in the manner of Auroux--Gayet--Mohnke \cite{AGM} and
Charest--Woodward \cite{CW}, and set $c_w(u)=\frac1k\,u\cdot D_w$.

\begin{definition}[singular symplectic complex; conditional on Problem~\ref{prob:relative}]
$C^{\rm par}_{\F[U,V]}(\mathcal H):=$ the Lagrangian Floer complex of $(L_\alpha,L_\beta)$ in $\mathcal M_{1/4}$, framed at a
point to remove reducibles (extended moduli spaces as in Manolescu--Woodward \cite{MW}), with
$\partial x=\sum\#\big(\cM(u)/\R\big)\,U^{c_w(u)}V^{c_z(u)}y$ over index-one strips.
\end{definition}

\begin{conjecture}[shape of the singular variant]\label{conj:symp}
\begin{enumerate}[label=(\alph*)]
\item $\partial^2=m_0(L_\beta)-m_0(L_\alpha)$, where $m_0(L)\in\F[U,V]$ counts Maslov-$2$ discs through a point weighted
by $U^{c_w}V^{c_z}$.  The complex is honest if these agree, and a matrix factorisation of
$W=m_0(L_\beta)-m_0(L_\alpha)$ otherwise.  (Compare Proposition~\ref{prop:d2}, where both equal $UV$.)
\item \textup{(No localisation to $Y$.)}  Inverting $V$ does \emph{not} return an invariant of $Y$.  The punctures are
built into $\mathcal M_\alpha$; they are not divisors in a space that knows only $Y$.  At $\alpha\to0$, $\mathcal M_\alpha$ is a
$\mathbb P^1\times\mathbb P^1$ flag bundle over the moduli of the underlying bundle, not that space with a
divisor removed.  The correct statement is a spectral sequence from the bivariant theory to $I^\natural(K)$
($U=V=1$).
\item \textup{(Symmetry)} The Hecke involution $\alpha\leftrightarrow\tfrac12-\alpha$ exchanges $\xi_w,\xi_z$ and hence $U,V$.
\end{enumerate}
\end{conjecture}

%----------------------------------------------------------------------
\subsection{The singular variant on the gauge side: two charges and a quadrant}\label{subsec:gauge-inst}
%----------------------------------------------------------------------
For singular connections on $\R\times(Y,K)$ with holonomy parameter $\tfrac14$, homotopy classes of paths between
lifted critical points form a torsor over $\Z^2\ni(k,l)$, the instanton and monopole numbers \cite{KM-knothom}.  The
index changes by $8k+4l$.  The topological energy at parameter $t\in[0,\tfrac12]$ of the configuration shifted by the
singular model, $B\mapsto B+(t-\tfrac14)\eta$, is $\kappa_t=k+2tl$ in Kronheimer--Mrowka's normalisation.  It is affine in
$t$, since the quadratic term $\int\mathrm{tr}(\eta\wedge d\eta)$ does not depend on $B$ and cancels in differences.  Set
\[
c_w:=\kappa_0=k,\qquad c_z:=\kappa_{1/2}=k+l,\qquad \kappa_{1/4}=\tfrac12(c_w+c_z),\qquad \mathrm{ind}=4(c_w+c_z).
\]
$(k,l)\mapsto(c_w,c_z)$ is an isomorphism $\Z^2\to\Z^2$.  The deck ring is $\F[U^{\pm1},V^{\pm1}]$, with $U=(k,l)=(1,-1)$ and
$V=(0,1)$: both have index $4$ and energy $\tfrac12$, and $T=V/U$ has index $0$ and energy $0$.  This is Daemi--Scaduto's
universal ring $\Z[U^{\pm1},T^{\pm1}]$ \cite{DS-equiv} in different coordinates.  What is new is the pair of real
functionals
\[
F_w:=\mathrm{CS}_0,\qquad F_z:=\mathrm{CS}_{1/2}\qquad\text{on the $\Z^2$-cover of }\cB(Y,K),
\]
whose differences along a trajectory are $c_w$ and $c_z$ up to the endpoint values.  Deck transformations translate
$(F_w,F_z)$ by $\Z^2$.  Every critical point therefore has a unique \emph{canonical lift} $\tilde a$ with
$(F_w,F_z)(\tilde a)\in[0,1)^2$.

\begin{lemma}\label{lem:quadrant}
If the differential of the universal complex (framed, with coefficients in $\F[U^{\pm1},V^{\pm1}]$) does not increase
$F_w$ or $F_z$ along any contributing trajectory, then the $\F[U,V]$-span of the canonical lifts is a subcomplex.  It is
free over $\F[U,V]$, and its localisation is the universal complex.
\end{lemma}
\begin{proof}
A trajectory from $\tilde a$ ends at $U^pV^q\tilde b$, with $F_w(U^pV^q\tilde b)=F_w(\tilde b)-p\le F_w(\tilde a)<1$, so $p>F_w(\tilde b)-1\ge-1$.
Likewise $q\ge0$.
\end{proof}

\begin{conjecture}[quadrant positivity]\label{conj:quadrant}
For metrics stretched along a doubly pointed Heegaard surface ($T\ge T(E)$, energy $\le E$), every index-$\le1$ singular
trajectory satisfies $c_w\ge0$ and $c_z\ge0$.
\end{conjecture}
In the adiabatic limit a trajectory becomes a strip in $\mathcal M_{1/4}$ together with bubbles.
\begin{itemize}[leftmargin=2em]
\item \emph{Strips.} $c_w,c_z\to\xi_w(u),\xi_z(u)$.  These are non-negative on spheres by Proposition~\ref{prop:nef}, and on
strips modulo Problem~\ref{prob:relative}.
\item \emph{Uhlenbeck bubbles} away from $\R\times K$ have $(k,l)=(m,0)$, so $(c_w,c_z)=(m,m)$: $UV$-type.
\item \emph{Singular bubbles} on $(\R^4,\R^2)$, conformally $(S^4,S^2)$.  Positivity is equivalent to
$|l|\le2\kappa_{1/4}$.  The minimal ones, $(k,l)=(0,1)$ and $(1,-1)$ with $\kappa=\tfrac12$, are exactly the pure-$V$ and pure-$U$
classes.  Heuristic proof: identify $(\C^2,\C\times0)$ with the complement of a line $\ell_\infty$ in
$(\mathbb P^2,\text{line})$.  Singular instantons framed at infinity should correspond to framed parabolic bundles
(Biquard \cite{Biquard}).  Then $\kappa_0=c_2(E)$ and $\kappa_{1/2}=c_2(E')$ for the elementary transform $E'$ along the line, and both are
$\ge0$ for framed sheaves (they are dimensions of ADHM data).  One must check that the elementary transform is
compatible with the framing at $\ell_\infty\cap\text{line}$.
\end{itemize}

\paragraph{An instanton $\Upsilon$.}
For $t\in[0,2]$ put $F_t:=(1-\frac t2)F_w+\frac t2F_z=\mathrm{CS}_{t/4}$.  If Conjecture~\ref{conj:quadrant} holds, every $F_t$ is a
filtration, and the spectral invariant of the distinguished (special-cycle) class defines $\Upsilon^\natural_K(t)$.  It
satisfies $\Upsilon^\natural_K(t)=\Upsilon^\natural_K(2-t)$ by the Hecke symmetry.  At $t=1$ it is a Chern--Simons-filtered invariant of
Daemi--Scaduto and Daemi--Imori--Sato--Scaduto--Taniguchi type \cite{DS-clasp,DISST}.  The other values of $t$ are the
content of positivity.  In this variant the structure belongs to the Khovanov/branched-cover side of the dictionary.
It is the singular analogue of the $t$-family, not of $\CFK$ itself.


%======================================================================
\section{Analytic details}\label{sec:analysis}
%======================================================================
The proposal reduces to analytic statements of four kinds:
\begin{enumerate}[label=(\alph*)]
\item an adiabatic limit with a knot transverse to the neck;
\item positivity and admissibility of the charges in that limit;
\item an intrinsic, neck-free description of the charges near $K$;
\item reducibles, compactness and invariance with the new weights.
\end{enumerate}
We take them in turn and say what is standard, what is new, and where the difficulty lies.

\subsection{The adiabatic limit with a transverse knot}
\emph{Monopoles.}  Use the metric $dt^2+ds^2+\varepsilon^2g_\Sigma$ on $\R\times N_T$, with an area-form perturbation of strength
$\tau$ satisfying the Bradlow condition for degree-$g$ vortices.  As $\varepsilon\to0$, solutions on the neck converge, away from
a finite bubbling set, to holomorphic maps $\R\times[-T,T]\to\mathrm{Sym}^g\Sigma$, the space of degree-$g$ $\tau$-vortices.  This is
Salamon's analysis for mapping tori \cite{Salamon-SW}, transplanted to a neck with Lagrangian ends.  What is new is
\emph{matching at the ends of the neck}.  One must show that the boundary values at $s=\pm T$ of finite-energy solutions
on the handlebody pieces $U_\alpha^\circ,U_\beta^\circ$ (cylindrical ends $[0,\infty)\times\Sigma$) converge to $\mathbb T_\alpha,\mathbb T_\beta$.  The model
statement is: \emph{the vortex boundary values of $\tau$-perturbed SW solutions on a handlebody with product end are the
divisors supported on the $\alpha$-curves, for a suitable choice of $\alpha$-curves determined by the metric}.  This is
the SW counterpart of ``flat connections on a handlebody form $L_\alpha$''.  It is the hardest single input and amounts
to the Atiyah--Floer conjecture for Seiberg--Witten.  The knot adds nothing to it: away from $s=\pm T$ the knot is the
pair of lines $\{w,z\}\times[-T,T]$, which are complex curves for the product structure.

\emph{Instantons.}  The analysis here is more advanced: Dostoglou--Salamon for mapping tori \cite{DoSa}, Salamon--Wehrheim for
Lagrangian boundary conditions \cite{SW-lag}, Daemi--Fukaya--Lipyanskiy for the $SO(3)$ Atiyah--Floer conjecture
\cite{DFL}.  The new ingredient is the determinant-line section along the strips, which is a smooth function of the
connection on the slices and so passes to the limit with the connection.  In its four-dimensional form, $c_z$ is the relative index of the Dirac
operator on $\R\times Y$ coupled to $\mathbb A\otimes\lambda_z$, with APS conditions at the ends.  It is therefore defined, additive and
locally constant for \emph{every} configuration.  This is cleaner than the monopole case: no admissibility is needed, and
only positivity remains, via Problem~\ref{prob:avoid}.  The analytic input is a gluing/adiabatic formula for this relative
index under neck-stretching, identifying it with the relative degree of $\mathcal D_z$ on the limiting strip plus bubble
contributions.  This is a families-index computation of Bismut--Freed/Dai type.

\subsection{Positivity in the limit}
\begin{lemma}[stability of non-negative zero counts]\label{lem:posvortex}
Let $\Omega=\R\times[-T,T]$, and let $s_\varepsilon$ be sections of a trivialised line bundle over $\Omega$ that converge in $C^0_{\rm loc}$ to a
holomorphic section $s_0$, uniformly on $\R\times\{\pm T\}$ and on $\{|t|\ge R_0\}\times[-T,T]$, where $|s_0|\ge c>0$.  Then the
signed zero count of $s_\varepsilon$ equals that of $s_0$ for $\varepsilon$ small, and it is $\ge0$.
\end{lemma}
\begin{proof}
On $\partial([-R_0,R_0]\times[-T,T])$ we have $|s_\varepsilon-s_0|<c\le|s_0|$ for small $\varepsilon$, so the winding numbers agree (Rouché).  The
zeros of $s_0$ are isolated and positive.
\end{proof}
The analytic work is in the hypotheses.  The first, $C^0$ convergence of the vortex component $\phi^-_z$ \emph{at the
specific point} $z$, needs elliptic estimates along $\Sigma$ that are uniform in $\varepsilon$.  These are available away from
bubbling points.  At a bubbling point lying on $\R\times\gamma^T_z$ one needs the bubble's own zero count.  In
$\mathrm{Sym}^g\Sigma$, sphere bubbles meet $V_w$ and $V_z$ equally and positively, and boundary degenerations carry $UV$.  The
gauge-theoretic statement to prove is that the energy-concentration profiles in the $\varepsilon\to0$ limit are vortex bubbles
of this type, and never configurations in which $\phi^-_z$ winds negatively.  The second hypothesis, uniformity in
$t$, is exponential convergence to critical points with constants independent of $\varepsilon$.  Since the
neck-stretched critical points correspond to transverse intersections of $\mathbb T_\alpha,\mathbb T_\beta$, this is a spectral-gap statement
for the Hessian, uniform in $\varepsilon$, of Dostoglou--Salamon type.

With Lemma~\ref{lem:bootstrap} these give Conjecture~\ref{conj:Mpos} for $T\ge T(E)$.  The energy window is essential:
there is no reason for positivity to hold uniformly over all energies at a fixed neck length.  The invariant is defined
as a direct limit (\S\ref{subsec:inv}).

\subsection{The intrinsic picture near $K$: critical holonomy and creases}\label{subsec:edge}
The neck is a scaffold.  For an invariant of $(Y,K)$ one wants to know what the charges are near $K$ when no splitting
is present.  The neck description says that along $\gamma_w$ the relevant component is $E_+(\dot K)$, and along $\gamma_z$ it is
$E_-(\dot K)$, switching at $p_0,p_3$.  Here is a local model in which this switching is built into the equations.

\paragraph{Aharonov--Bohm holonomy.}  Couple the spin$^c$ connection to a flat $U(1)$ connection on $Y\setminus K$ with
meridional holonomy $e^{2\pi i\beta}$.  For $U(1)$ this is the limit of an \emph{exact} perturbation, since the Thom form of a
null-homologous $K$ is exact.  So $\HM$ is unchanged for smooth smearings, which is \ref{Q:base}.  In the singular limit,
transversally to $K$ the Dirac operator is the Aharonov--Bohm operator on $\C$ with flux $\beta\in(0,1)$.  Its local
solutions near $r=0$ have leading terms
\[
\Phi\sim c_+\,r^{-\beta}\,\mathbf e_+ +c_-\,r^{\beta-1}\,\mathbf e_-+O(r^{\min(1-\beta,\beta)}),\qquad \mathbf e_\pm\in E_\pm(\dot K),
\]
and both singular terms are in $L^2$.  The operator therefore has a family of closed extensions, parametrised by
Lagrangians in the space of edge data $(c_+,c_-)\in\Gamma(K;E_+\oplus E_-)$, with the symplectic form coming from Green's formula,
which pairs $c_+$ with $c_-$.  Two extensions are distinguished: $\mathcal E_w=\{c_-=0\}$, which tends to the smooth theory as
$\beta\to0$, and $\mathcal E_z=\{c_+=0\}$, which tends to the smooth theory after the singular gauge transformation $e^{i\theta}$ as
$\beta\to1$.  At the \emph{critical holonomy} $\beta=\tfrac12$ both singular terms are $r^{-1/2}$, and the two extensions are
exchanged by the reflection $\theta\mapsto-\theta$.

\paragraph{Creases.}  The proposal is to use $\beta=\tfrac12$ with the boundary condition $\mathcal E_w$ along $\gamma_w$ and $\mathcal E_z$ along
$\gamma_z$.  The condition switches at the two \emph{creases} $p_0,p_3$, and the surviving edge coefficient is the section
$c=(c_+|_{\gamma_w},c_-|_{\gamma_z})$.  This is precisely the neck description, now intrinsic, and the charges become
\[
n_w=\#\,c_+^{-1}(0)\cap(\R\times\gamma_w),\qquad n_z=\#\,c_-^{-1}(0)\cap(\R\times\gamma_z).
\]
Their difference is a topological count for the glued line bundle over $\R\times K$, which gives \ref{M:alex}.  Positivity on
both arcs is what prevents zeros from crossing a crease: a positive zero of $c_+$ on the $\gamma_w$ side would arrive as a
negative zero of $c_-$ on the $\gamma_z$ side.  So the creases behave like the Lagrangian walls of the cylindrical
reformulation of Heegaard Floer.  On a cobordism with surface $F$, the boundary condition can switch along curves on
$F$.  These \emph{crease curves are Zemke's dividing arcs} \ref{M:cob}, and the index contributions of the
$\mathcal E_w$- and $\mathcal E_z$-regions should reproduce the $\tilde\chi(F_w)$ and $\tilde\chi(F_z)$ terms of his grading formula.

\paragraph{What fails, and what must be proved.}
\begin{enumerate}[label=(\roman*),leftmargin=2em]
\item \emph{Edge calculus.}  The linearised operators are edge operators in Mazzeo's sense \cite{Mazzeo} along $K$ (in $3$d) and
$\R\times K$ (in $4$d), with indicial roots $\pm\frac12+\Z$ for the Dirac part at $\beta=\frac12$.  Fredholmness on weighted spaces with
the mixed boundary condition, and an index formula, are standard in principle.  The crease is a codimension-$3$
stratum in $4$d where the boundary condition jumps.  Its local model is a two-dimensional problem with corner-type
indicial roots.  It must be shown that the index jumps by an amount independent of the configuration, which is the
analogue of the fixed Maslov contribution of a basepoint region.
\item \emph{$C^0$ bound.}  The standard bound $|\Phi|^2\le\max(-s_g,0)$ fails, since $|\Phi|\sim r^{-1/2}$ near $K$.  Replace it by a
weighted maximum principle for $r|\Phi|^2$, using the Weitzenböck formula on $Y\setminus K$ together with the boundary term from
Green's formula.  That boundary term is $\mathrm{Re}\int_K\langle c_+,c_-\rangle$ up to constants, and it \emph{vanishes} for $\mathcal E_w$ and for
$\mathcal E_z$.
\item \emph{Edge bubbling.}  The SW equations are scale-invariant on $\R^4$, and only the $C^0$ bound prevents bubbling.
Near $K$, rescaled limits are finite-energy solutions on $\R^4$ with Aharonov--Bohm flux $\frac12$ along $\R^2$.  For a
one-sided condition the Weitzenböck identity $\int|\nabla\Phi|^2+\frac14|\Phi|^4=(\text{edge term})=0$ makes them trivial.  At a crease
the edge term need not vanish, so nontrivial \emph{crease bubbles} may exist.  The prediction is that they exist and
carry weight $UV$, cancel in pairs, and are the gauge-theoretic boundary degenerations of
Proposition~\ref{prop:d2}.  Classifying finite-energy solutions of this local model is a concrete,
finite-dimensional-symmetry problem.  By $\mathrm{SO}(2)\times\mathrm{SO}(2)$-invariance it reduces to ODE.
\item \emph{Positivity without a neck.}  Unlike in the neck limit, the edge coefficients satisfy no local
$\bar\partial$-equation: leading edge data are free boundary data.  Intrinsic positivity of $n_w,n_z$ therefore requires a
global argument, for instance a Kähler-type model metric on a tubular neighbourhood of $\R\times K$ combined with a
large perturbation localised there, in Taubes' style.  We regard the neck construction as primary and this as its
eventual intrinsic replacement.
\end{enumerate}

\paragraph{Bordered consistency check.}  If one cuts along $T^2=\partial\nu K$ instead of $\Sigma$, the boundary phase space is the
dual torus $\Tdual=H^1(T^2;\R/\Z)$ of flat $U(1)$ connections.  The flat-torus Dirac operator has kernel at exactly one point
$\theta_0\in\Tdual$ for a fixed spin structure, and there the kernel is $\C\oplus\C$, one constant spinor of each chirality.  Three
facts can be checked directly:
\begin{enumerate}[label=(\arabic*)]
\item the solid-torus Lagrangian avoids $\theta_0$;
\item the reducible locus of the knot complement, $\{\mathrm{hol}_\lambda=1\}$, avoids $\theta_0$;
\item the Aharonov--Bohm filling at $\beta=\frac12$ passes through $\theta_0$.
\end{enumerate}
This matches the Hanselman--Rasmussen--Watson picture.  Fillings miss the marked point, $\widehat{\HFK}$ is computed by a meridian
through it, the essential component of the curve is the reducible line, and the two chiral kernels at $\theta_0$ are the two
sides of the marking that carry $w$ and $z$.  It is an independent check that critical holonomy is the right local
model.

\subsection{Reducibles, compactness and invariance}\label{subsec:inv}
\begin{itemize}[leftmargin=2em]
\item \emph{Blow-up.}  In Kronheimer--Mrowka's $\sigma$-model the charges are defined with the normalised spinor $\phi$, so they
make sense on reducible trajectories.  For a flow from eigenvector $\psi_\lambda$ to eigenvector $\psi_\mu$, $n_z$ is the difference of the
windings of $\pi_{E_-}\psi_\lambda$ and $\pi_{E_-}\psi_\mu$ along $\gamma_z^T$, or of the edge coefficients in the intrinsic model.  Genericity of
eigenvectors, meaning non-vanishing of a section of a real rank-$2$ bundle along an arc, is a codimension-one condition
in $1$-parameter families.  It must be arranged by perturbations that preserve the structure near $K$.
\item \emph{Energy windows and direct limits.}  Take $C_{T,E}$ to be the subquotient generated in energy $\le E$.  For
$T'\ge T\ge T(E')$ the continuation maps $C_{T,E}\to C_{T',E'}$ are $V$-weighted and positive, and the invariant is
$\varinjlim_E H_*(C_{T(E),E})$.  This copies Taubes' passage from $\HM$ to ECH via action-filtered complexes.  The point to
check is that continuation trajectories also satisfy positivity once $T$ is large relative to the energy.  That is the
1-parameter version of Lemma~\ref{lem:posvortex}.
\item \emph{Splitting independence.}  Stabilisations and handleslides away from $w,z$ are realised by stretching along
the extra surfaces of a Heegaard triple.  The corresponding triangle counts are $V$-weighted by the same charges on the
triangle's neck strips.  Moves across $z$ are exactly those along which admissibility fails, as in Heegaard Floer.
\item \emph{Orientations.}  Work over $\F$ first.  Over $\Z$, Kronheimer--Mrowka's orientation of the determinant line is
unaffected by the weights, but crease bubbles (if they exist) must be shown to cancel with signs.  In Heegaard Floer
the two boundary degenerations do cancel with signs.
\end{itemize}

%======================================================================
\section{Tests, first cases, and a roadmap}\label{sec:tests}
%======================================================================

\subsection{Monopole versus instanton}
\begin{center}\small
\begin{tabular}{p{3.3cm}p{5.2cm}p{5.2cm}}
\toprule
& monopole model (\S\ref{subsec:monopole}) & instanton model (\S\ref{subsec:instanton})\\
\midrule
base theory & $\widehat{CM}(Y)$, $U$ from $S^1$-equivariance & $SO(3)$-equivariant $I^-(Y)$, $U$ of degree $4$\\
$z$-charge & zeros of the vortex component $\pi_{E_-(\dot K)}\phi$ on $\R\times\gamma_z^T$ & relative index of Dirac coupled to $\mathbb A\otimes\lambda_z$\\
always defined? & no: admissibility needed & yes\\
Atiyah--Floer target & $\mathrm{Sym}^g\Sigma$, $V_z$ (this \emph{is} $\CFK$) & extended $\mathcal R(\Sigma)$, theta-type $D_z$\\
how $K$ enters & handlebody metrics select $\alpha,\beta$; position of $z$ & trivialisations of $\mathcal D_z$ via the arc strings\\
positivity from & vortex holomorphicity (Lemma~\ref{lem:posvortex}) & compatible divisors (Problem~\ref{prob:avoid})\\
extra subtlety & reducible tower carries $V$-weights & deck (deg $8$) vs.\ equivariant (deg $4$) variable\\
\bottomrule
\end{tabular}
\end{center}

\subsection{Consistency checks a construction must pass}
\begin{enumerate}[label=(C\arabic*),leftmargin=2.6em]
\item \emph{Hat level.}  $H_*\big(\CFK^{\HM}\otimes\F[U,V]/(U,V)\big)\cong KHM(Y,K)$ and
$H_*\big(\CFK^{I}\otimes\F[U,V]/(U,V)\big)\cong KHI(Y,K)$.  The first is a known group, isomorphic to $\widehat{\HFK}$
\cite{Lekili,BS-equiv,CGH-sutured}.  A direct gauge-theoretic proof would compare the $V=0$, framed theory on $g_T$ with the
sutured closure by excision along the tori $\partial\nu\gamma_z$.  This is the most accessible nontrivial theorem.
\item \emph{$V=1$.}  Built in (Theorem~\ref{thm:Mstructure}(b)).
\item \emph{$U$ inverted.}  Must give $\HM(Y)$ seen from $z$.  This is a statement about edge windings of Dirac
eigenvectors (\S\ref{subsec:monopole}), and it is genuinely nontrivial.
\item \emph{Gradings.}  $\gr_w$ is Kronheimer--Mrowka's absolute $\mathbb Q$-grading, and $\gr_z=\gr_w-2A$.  The Alexander function
must agree with $\frac12\langle c_1,\hat S\rangle$ for relative spin$^c$ structures; in the instanton model, with the generalised
eigenvalues of $\mu(\hat S)$ on $KHI$.
\item \emph{Reversal and flip.}  Reversing $K$ exchanges $\gamma_w,\gamma_z$ and $E_\pm(\dot K)$, hence $U\leftrightarrow V$.  Moving the creases
$p_0,p_3$ once around $K$ should give a monodromy $\varsigma\simeq1+\Phi\Psi$, where $\Phi$ and $\Psi$ count trajectories weighted by
$n_w$ and $n_z$.  This is Sarkar's formula, now as the monodromy of a family of gauge-theoretic data.
\item \emph{Surgery.}  The $2$-handle cobordism $W_N$ with its core surface carries $V$-weighted maps.  The large surgery
statement $\widehat{\HM}(Y_N(K),[s])\cong H_*(A^-_s)$ for $N\gg0$ is a cobordism-map computation in which the core disc is cut by
the arcs $\gamma_w,\gamma_z$, so that its two halves are Zemke's $F_w,F_z$.  The mapping cone then needs only the flip, (C5).
\end{enumerate}

\subsection{First cases}
\paragraph{Unknot, round $S^3$.}  Put $K=\{z_2=0\}\subset S^3\subset\C^2$.  The torus $T^2$ acting on $(z_1,z_2)$ preserves $K$, the round
metric, the reducible, and (in the intrinsic model of \S\ref{subsec:edge}) the Aharonov--Bohm data.  Dirac eigenvectors
are $T^2$-weight vectors $e^{i(m_1\varphi_1+m_2\varphi_2)}f(\rho)$, with $z_1=\cos\rho\,e^{i\varphi_1}$, $z_2=\sin\rho\,e^{i\varphi_2}$, and the
eigenvalue problem separates into ODEs in $\rho$.  Along $K=\{\rho=0\}$ the winding of an eigenvector is its longitudinal
weight $m_1$, up to a framing shift.  The edge exponents are governed by the meridional weight $m_2\in\Z+\tfrac12$ at critical
holonomy.  The prediction $\CFK=\F[U,V]$ becomes an explicit statement about the joint ordering of eigenvalues and
weights $(m_1,m_2)$.  This is the cleanest first computation, and it tests (C3), (C4) and the critical-holonomy
model at once.

\paragraph{Genus one: $(1,1)$-knots.}  Every $(1,1)$-knot, in particular every torus knot in $S^3$, has a genus-one doubly
pointed diagram.
\begin{itemize}[leftmargin=2em]
\item \emph{Monopole.}  The vortex moduli space is $\mathrm{Sym}^1T^2=T^2$, the strips are immersed bigons, and the neck is
$\R\times[-T,T]\times T^2$, flat.  The required adiabatic analysis is closest to the product formulae along tori of
Morgan--Mrowka--Szab\'o \cite{MMS} and Taubes.
\item \emph{Instanton.}  $\mathcal R(T^2)$ is the pillowcase, $L_\alpha,L_\beta$ are arcs between corners (the reducibles, resolved by
framing or holonomy perturbation as in Hedden--Herald--Kirk \cite{HHK1}), and $\mathcal D_p$ has a single compatible zero.  So
$D_w,D_z$ are \emph{two marked points of the pillowcase}, separated by the image of $w-z$ under $\Sigma\to\mathrm{Jac}(\Sigma)$.  The knot is
recorded by the isotopy class of the pair of marked points relative to the arcs.  The conjecture is that
\[
\CFK^{I}_{\F[U,V]}(K)\simeq CF\big(L_\alpha,L_\beta;\ \text{pillowcase},\ U^{\#D_w}V^{\#D_z}\big),
\]
which is the instanton counterpart of the genus-one combinatorics of $\CFK$.  Its $U=V=0$ part should reproduce
Li--Ye's computation of $KHI$ for $(1,1)$-knots \cite{LY-heegaard}.  This is an explicit, finite check, and the genus-one
case is the natural place to make the proposal rigorous first.
\end{itemize}

\paragraph{Trefoil.}  $\CFK(T_{2,3})=\langle a,b,c\rangle$ with $\partial b=Ua+Vc$.  In the pillowcase model there must be exactly two bigons
from $b$: one covering $D_w$ once and one covering $D_z$ once.  In the monopole model, one trajectory from $b$ has $n_z=1$ and
the other has $n_z=0$, with the $U$ coming from the reducible tower.

\subsection{Roadmap}
\begin{enumerate}[leftmargin=2em]
\item Algebra: Theorem~\ref{thm:Mstructure} and its instanton analogue, assuming positivity.  This is routine.
\item Instanton charges: define $c_w,c_z$ as relative Dirac indices for all configurations, prove additivity, and
compute them on bubbles ($UV$-type).  This needs index theory but no adiabatic limit.
\item Genus one, instanton: solve Problem~\ref{prob:avoid} on the pillowcase, which is a $2$-dimensional problem, prove
positivity by the adiabatic limit for $T^2$-necks, and verify (C1) for torus knots.
\item Genus one, monopole: the same with $\mathrm{Sym}^1T^2=T^2$.  Compare with $\widehat{\HFK}$ via $KHM$.
\item The intrinsic critical-holonomy model: edge calculus, the weighted $C^0$ bound, and classification of crease
bubbles (\S\ref{subsec:edge}(i)--(iii)).  Prove that the neck and intrinsic charges agree.
\item Higher genus: the general Atiyah--Floer comparison (Conjecture~\ref{conj:MAF}), large surgery, flip, and
mapping cone.
\end{enumerate}

\subsection{What is new here, in one paragraph}
On the Heegaard Floer side, $\CFK$ is $\CF^-(Y)$ together with an Alexander function and a one-sided positivity
(Lemma~\ref{lem:reform}).  On the gauge side, the proposal supplies these as follows.
\begin{itemize}[leftmargin=2em]
\item \emph{Charges.}  The $z$-charge is a topological count of vortex zeros (monopoles) or a relative Dirac index with a
unit flux string along the $z$-arc of $K$ (instantons).  The Alexander function is the source-free difference of the
$w$- and $z$-charges.
\item \emph{Positivity.}  It comes from the Atiyah--Floer limit, by holomorphicity together with integrality.
\item \emph{Arcs and creases.}  The two arcs of $K$ meet at creases, which are the gauge-theoretic home of the Heegaard
Floer basepoint walls and of Zemke's dividing sets.
\item \emph{Local model at $K$.}  Critical Aharonov--Bohm holonomy $\beta=\frac12$, whose two edge modes are the two chiral
components that $w$ and $z$ see.
\item \emph{Singular variant.}  It gives a separate, Khovanov-side two-variable theory, in which positivity is nefness of
the two ends of the parabolic Kähler segment and $t=4\alpha$.
\end{itemize}

\begin{thebibliography}{99}\small

\bibitem{OS-knots} P.~Ozsv\'ath, Z.~Szab\'o, \emph{Holomorphic disks and knot invariants}, Adv.\ Math.\ 186 (2004).
\bibitem{Rasmussen} J.~Rasmussen, \emph{Floer homology and knot complements}, PhD thesis, Harvard, 2003.
\bibitem{Zemke-link} I.~Zemke, \emph{Link cobordisms and functoriality in link Floer homology}, arXiv:1610.05207, J.\ Topol.
\bibitem{Zemke-grading} I.~Zemke, \emph{Link cobordisms and absolute gradings on link Floer homology}, Quantum Topol.\ 10 (2019); arXiv:1701.03454.
\bibitem{Sarkar} S.~Sarkar, \emph{Moving basepoints and the induced automorphisms of link Floer homology}, AGT 15 (2015).
\bibitem{HM-inv} K.~Hendricks, C.~Manolescu, \emph{Involutive Heegaard Floer homology}, Duke Math.\ J.\ 166 (2017).
\bibitem{JTZ} A.~Juh\'asz, D.~Thurston, I.~Zemke, \emph{Naturality and mapping class groups in Heegaard Floer homology}, arXiv:1210.4996, Mem.\ AMS.
\bibitem{OSS} P.~Ozsv\'ath, A.~Stipsicz, Z.~Szab\'o, \emph{Concordance homomorphisms from knot Floer homology}, Adv.\ Math.\ 315 (2017).
\bibitem{Hom} J.~Hom, \emph{Bordered Heegaard Floer homology and the tau-invariant of cable knots}, J.\ Topol.\ 7 (2014).
\bibitem{DHST} I.~Dai, J.~Hom, M.~Stoffregen, L.~Truong, \emph{More concordance homomorphisms from knot Floer homology}, Geom.\ Topol.\ 25 (2021).
\bibitem{NiWu} Y.~Ni, Z.~Wu, \emph{Cosmetic surgeries on knots in $S^3$}, J.\ reine angew.\ Math.\ 706 (2015).
\bibitem{HRW} J.~Hanselman, J.~Rasmussen, L.~Watson, \emph{Bordered Floer homology for manifolds with torus boundary via immersed curves}, J.\ Amer.\ Math.\ Soc.\ 37 (2024).
\bibitem{Hanselman} J.~Hanselman, \emph{Knot Floer homology as immersed curves}, arXiv:2305.16271.

\bibitem{KM-book} P.~Kronheimer, T.~Mrowka, \emph{Monopoles and three-manifolds}, Cambridge Univ.\ Press, 2007.
\bibitem{KM-sutures} P.~Kronheimer, T.~Mrowka, \emph{Knots, sutures and excision}, J.\ Differential Geom.\ 84 (2010).
\bibitem{KM-alex} P.~Kronheimer, T.~Mrowka, \emph{Instanton Floer homology and the Alexander polynomial}, AGT 10 (2010).
\bibitem{Lim} Y.~Lim, \emph{Instanton homology and the Alexander polynomial}, Proc.\ AMS 138 (2010).
\bibitem{KM-knothom} P.~Kronheimer, T.~Mrowka, \emph{Knot homology groups from instantons}, J.\ Topol.\ 4 (2011).
\bibitem{KM-unknot} P.~Kronheimer, T.~Mrowka, \emph{Khovanov homology is an unknot-detector}, Publ.\ Math.\ IH\'ES 113 (2011).
\bibitem{KM-conc} P.~Kronheimer, T.~Mrowka, \emph{Instantons and some concordance invariants of knots}, arXiv:1910.11129.
\bibitem{Lekili} Y.~Lekili, \emph{Heegaard Floer homology of broken fibrations over the circle}, Adv.\ Math.\ 244 (2013).
\bibitem{BS-equiv} J.~Baldwin, S.~Sivek, \emph{On the equivalence of contact invariants in sutured Floer homology theories}, Geom.\ Topol.\ 25 (2021).
\bibitem{CGH-sutured} V.~Colin, P.~Ghiggini, K.~Honda, \emph{Sutured Heegaard Floer and embedded contact homologies are isomorphic}, arXiv:2403.16492.
\bibitem{KLT} \c{C}.~Kutluhan, Y.-J.~Lee, C.~Taubes, \emph{HF=HM I--V}, Geom.\ Topol.\ 24 (2020).
\bibitem{Li-minus} Z.~Li, \emph{Knot homologies in monopole and instanton theories via sutures}, J.\ Symplectic Geom.\ 19 (2021).
\bibitem{GZ} S.~Ghosh, I.~Zemke, \emph{Connected sums and directed systems in knot Floer homologies}, arXiv:2303.06491.
\bibitem{DS-equiv} A.~Daemi, C.~Scaduto, \emph{Equivariant aspects of singular instanton Floer homology}, Geom.\ Topol.\ 28 (2024).
\bibitem{DS-clasp} A.~Daemi, C.~Scaduto, \emph{Chern--Simons functional, singular instantons, and the four-dimensional clasp number}, JEMS 26 (2024).
\bibitem{DISST} A.~Daemi, N.~Imori, K.~Sato, C.~Scaduto, M.~Taniguchi, \emph{Instantons, special cycles, and knot concordance}, Geom.\ Topol.\ 29 (2025).
\bibitem{ME} M.~Miller Eismeier, \emph{Equivariant instanton homology}, arXiv:1907.01091.
\bibitem{DME} A.~Daemi, M.~Miller Eismeier, \emph{Instantons and rational homology spheres}, arXiv:2210.14071.
\bibitem{BS-conc} J.~Baldwin, S.~Sivek, \emph{Framed instanton homology and concordance}, J.\ Topol.\ 14 (2021).
\bibitem{LY-heegaard} Z.~Li, F.~Ye, \emph{Instanton Floer homology, sutures, and Heegaard diagrams}, J.\ Topol.\ 15 (2022).
\bibitem{LY-large} Z.~Li, F.~Ye, \emph{SU(2) representations and a large surgery formula}, arXiv:2107.11005.
\bibitem{LY-surgery} Z.~Li, F.~Ye, \emph{Knot surgery formulae for instanton Floer homology I: the main theorem}, Geom.\ Topol.\ 29 (2025).
\bibitem{HHK1} M.~Hedden, C.~Herald, P.~Kirk, \emph{The pillowcase and perturbations of traceless representations of knot groups}, Geom.\ Topol.\ 18 (2014).
\bibitem{HHK2} M.~Hedden, C.~Herald, P.~Kirk, \emph{The pillowcase and traceless representations of knot groups II}, arXiv:1501.00028.
\bibitem{CHKK} G.~Cazassus, C.~Herald, P.~Kirk, A.~Kotelskiy, \emph{The correspondence induced on the pillowcase by the earring tangle}, arXiv:2010.04320.
\bibitem{Horton} H.~Horton, \emph{A symplectic instanton homology via traceless character varieties}, arXiv:1611.09927.
\bibitem{MW} C.~Manolescu, C.~Woodward, \emph{Floer homology on the extended moduli space}, in \emph{Perspectives in Analysis, Geometry, and Topology}, Progr.\ Math.\ 296 (2012).
\bibitem{Wang} D.~Wang, \emph{Monopole Floer homology for 3-manifolds with toroidal boundary}, arXiv:2004.06227; see also arXiv:2005.04333, arXiv:2010.04318.
\bibitem{SW-lag} D.~Salamon, K.~Wehrheim, \emph{Instanton Floer homology with Lagrangian boundary conditions}, Geom.\ Topol.\ 12 (2008).
\bibitem{DFL} A.~Daemi, K.~Fukaya, M.~Lipyanskiy, \emph{Lagrangians, SO(3)-instantons and the Atiyah--Floer conjecture}.
\bibitem{DoSa} S.~Dostoglou, D.~Salamon, \emph{Self-dual instantons and holomorphic curves}, Ann.\ of Math.\ 139 (1994).
\bibitem{Salamon-SW} D.~Salamon, \emph{Seiberg--Witten invariants of mapping tori, symplectic fixed points, and Lefschetz numbers}, Turkish J.\ Math.\ 23 (1999).
\bibitem{MS} V.~Mehta, C.~Seshadri, \emph{Moduli of vector bundles on curves with parabolic structures}, Math.\ Ann.\ 248 (1980).
\bibitem{Pauly} C.~Pauly, \emph{Espaces de modules de fibr\'es paraboliques et blocs conformes}, Duke Math.\ J.\ 84 (1996).
\bibitem{Biquard} O.~Biquard, \emph{Sur les fibr\'es paraboliques sur une surface complexe}, J.\ London Math.\ Soc.\ (2) 53 (1996).
\bibitem{AGM} D.~Auroux, D.~Gayet, J.-P.~Mohsen, \emph{Symplectic hypersurfaces in the complement of an isotropic submanifold}, Math.\ Ann.\ 321 (2001).
\bibitem{CW} F.~Charest, C.~Woodward, \emph{Floer cohomology and flips}, Mem.\ Amer.\ Math.\ Soc.\ 279 (2022), no.\ 1372; arXiv:1508.01573.
\bibitem{MMS} J.~Morgan, T.~Mrowka, Z.~Szab\'o, \emph{Product formulas along $T^3$ for Seiberg--Witten invariants}, Math.\ Res.\ Lett.\ 4 (1997).
\bibitem{Mazzeo} R.~Mazzeo, \emph{Elliptic theory of differential edge operators I}, Comm.\ PDE 16 (1991).

\end{thebibliography}

\end{document}
